\chapter{Background}
\label{chapter:Background}

In this chapter we will introduce the weak memory, the Lean 4 programming langauge and its metaprogramming framework, and the mathematical backgrounds needed in later chapters. First, we introduce the weak memory models by some basic examples and weak behivours, Section \ref{Background:LKMM}

\section{Weak Memory Model}
There was a widely held belief in the software community that we did not need to optimize software; we could focus on abstraction, and hardware improvements would deliver performance gains for us. This belief largely held during the single-threaded era, when increasing clock speeds and improved microarchitectures steadily accelerated programs. However, when single-core processors hit power and energy-efficiency limits, hardware designers shifted toward adding more cores to improve performance. At that point, synchronization and memory-consistency abstractions were not yet well defined for compiler and system engineers. Vendors relied heavily on experience and natural-language descriptions to explain hardware behavior. Because natural language is inherently ambiguous, it became difficult to reason precisely about many concurrent program behaviors.

The compiler will normally optimize the code aggressively for performance as long as the programmer can’t observe the difference between the unoptimized and optimized code, but in fact, the code is not equivalent. In a uniprocessor system it works fine, but when it comes to a multithreaded program the code can cause bugs. Also, the hardware itself will run the program out of order to optimize execution, or just speculatively execute. Optimizing compilers have great freedom in the way they translate source code to object code. They are allowed to apply transformations that add memory accesses, eliminate accesses, combine them, split them into pieces, or move them around.

Let’s consider the following example:
\[
\text{Initially: } X = 0,\; Y = 0
\]

\[
\begin{array}{c|c}
\text{Thread 0} & \text{Thread 1} \\
\hline
X := 1          & Y := 1 \\
a := Y          & b := X
\end{array}
\]
This is a Store-Buffering program in which the variables X and Y are initialized to 0. Each thread first assigns the value 1 to one variable and then reads the other variable. From the programmer’s perspective, one might expect both a and b to be 1, assuming the program executes in the order in which the operations are written. However, due to underlying hardware architectures and compiler optimizations, this assumption does not necessarily hold. It is possible for both a and b to be 0, and it is also possible for a and b to observe different values.

To help describe these weak memory behaviors, Jade proposed an axiomatic semantics called Cat, which uses binary relations. In the next section, we introduce this axiomatic semantics and its specification language (Cat) in section \ref{Background:AWMM}.

\subsection{Axiomatic Weak Memory Model}
\label{Background:AWMM}
Memory models are generally categorized into two classes: operational and axiomatic. Operational models are described using buffers and queues; they are abstractions derived from hardware behavior \citet{10.1145/1993498.1993520OperationalModel}. Axiomatic memory models, in contrast, describe the disallowed behaviors of program executions, usually by constraining various relations on memory accesses \citet{10.1145/2627752}.

The axiomatic model uses the Event as the primitive type, they use event to give sementics to each single instruction, the event contains the thread id, the address, the operation (Fence, read or write) and instruction value, etc. 

\subsection{Linux Kernel Memory Model}
\label{Background:LKMM}
The Linux system supports a large range of architectures, such as ARM, RISC-V, IBM Power, AMD/Intel, and each architecture has their own memory models and thus disagree on the values that a read can return in a multiprocessor system, it's very important to support all of them in an efficient and secure way, a bad memory model can lead to ineffiencian problem like [], and secure problems like []. The LKMM has to be fairly permissive, in the sense that any behavior allowed by one of these architectures also has to be allowed by the LKMM.

The Linux Kernel memory model (LKMM) is formalized using a Domain Specific Language (DSL) called Cat \citet{alglave2018frightening}, before that the Linux kernel developers use the experiences from maintainers mostly, and they found out they need to find a way that can guide their develpoment. 
And they define the Linux kernel memory models are defined in .cat and .bell files in this link.

\subsubsection{Observation, code equivalance, and bugs}



Let's see this example in Pseudo Code (Message Passing):

buf = 0, flag = 0;

	P0()
	{
		write(buf, 1);
		write(flag, 1);
	}

	P1()
	{
		int r1;
		int r2 = 0;

		r1 = load(flag);
		if (r1)
			r2 = load(buf);
	}

Normally, we expect that in P1 the values of r1 and r2 are always consistent; if they are different, we find a bug. But this assumption can be broken easily by compiler and hardware optimizations.

First, the compiler may optimize this because these two write operations are independent, so the generated assembly code may write the flag first and then write the buf, because from the compiler’s perspective, swapping them has no observable effects. But if we just let the compiler optimize it, in P1 we may read buf as 0 after reading flag as 1.

Second, even if we don’t let the compiler do any optimizations, the hardware will still break the code. In some weak memory model architectures like ARM and RISC-V, they may reorder the writes, which will cause the same problem as the compiler optimization. But the TSO model will prevent this, because the TSO model let other threads observe the writes in program order (The abstract model is that they put all the writes in the queue, so there is no reordering of write-write (WW) operations). So, the outcome highly depends on the hardware the code executes on.

The most foundamental way to access the memory in Linux kernel is READ\_ONCE and WRITE\_ONCE, as they prevent the optimization of compiler to split or combine the memory access operation, for example, the compiler may combine many read opreations at once to make the use of the memory bandwidth, the ONCE is used to prevent such optimizations.

\section{Mathematical Backgrounds}

The formalization used in this thesis is based on \emph{binary relations}. Relationships between elements of two sets are represented using a structure called a binary relation. There are two common methods for representing relations: using square matrices and using directed graphs. In this thesis, we mainly use the second method, i.e., directed graphs, to represent binary relations.

We use the notation \(a \, R \, b\) to denote that \((a,b) \in R\), and \(a \not R b\) to denote that \((a,b) \notin R\).

In this thesis, the elements \(a\) and \(b\) are both \emph{events}, and these binary relations compose the directed execution graph, as described in Section~\ref{Background:AWMM}. Two relations can be composed together to obtain a new directed graph. Relation composition plays an important role in the \texttt{cat} language. The definition of the composition of binary relations is as follows:

\begin{definition}
Let \(R\) be a relation from a set \(A\) to a set \(B\), and let \(S\) be a relation from \(B\) to a set \(C\).
The composite of \(R\) and \(S\) is the relation consisting of ordered pairs \((a,c)\), where
\(a \in A\), \(c \in C\), and there exists an element \(b \in B\) such that \((a,b) \in R\) and
\((b,c) \in S\). We denote the composite of \(R\) and \(S\) by \(S ; R\).
\end{definition}

Computing the composite of two relations requires finding elements that are the second component of ordered pairs in the first relation and the first component of ordered pairs in the second relation. This property is widely used in many memory models. For example, there is a relation called \emph{fr} (from-read), which represents that a read instruction observes a value that is later overwritten by another write instruction, the write instructions to the same address are totally ordered by the coherence order.

In \texttt{cat}, we define
\[
fr = co ; rf^{-1},
\]
where \(rf^{-1}\) denotes the inverse of the relation \(rf\), meaning that the order of each pair is reversed. Intuitively, this composition connects a read event to a write event that comes later in coherence order than the write from which the read obtained its value.

We provide an illustration of this construction in the following figure:

\emph{[provide a transition of connection of graphs]}

After composing different relations, we can obtain a complete execution graph. As we said in Section~\ref{Background:AWMM}, we need to filter out disallowed executions using a set of judgments. cat provides three primitive judgments: Acyclic, Irreflexive, and Empty. Among them, Acyclic is the most important, because our memory models are described by binary relations that compose the execution graph. A cycle in this graph usually means that some data that should have been read first was overwritten by later operations.

Models defined using binary relations can be viewed as directed graphs. When we illustrate them in a diagram, cycles are easy to identify by visually connecting the points. However, to represent this structure using binary relations, we need a formal way to express both direct connections and paths that return to earlier points. For this purpose, we rely on two properties: \emph{reflexivity} and \emph{transitive closure}.

\paragraph{Reflexivity.}
A relation \( R \) on a set \( X \) is said to be \emph{reflexive} if, for every \( x \in X \), \( (x, x) \in R \); that is,
\[
\forall x \in X,\ x R x. 
\]

\paragraph{Transitivity.}
A binary relation \( R \) on a set \( X \) is \emph{transitive} if, for all elements \( a, b, c \in X \), whenever \( (a, b) \in R \) and \( (b, c) \in R \), then \( (a, c) \in R \); formally,
\[
\forall a, b, c \in X,\ a R b \land b R c \;\Rightarrow\; a R c.
\]

A closure is 


\subsubsection{Comparison between models in different architectures}

From a high-level perspective, when comparing different models, we check whether one model allows more program behaviors than another. If it does, we say that this model is weaker than the other one. A program generates a set of candidate executions, denoted by \( S^c \), and the memory model filters these to determine which executions are allowed. A weaker memory model should include all executions allowed by a stricter memory model. Formally, we require

\[
\{ X \mid X \in S^c \text{ and } \text{allowed}_{\text{strict}}(X) \}
\;\subseteq\;
\{ X \mid X \in S^c \text{ and } \text{allowed}_{\text{weak}}(X) \}.
\]

This captures the idea that every execution permitted by the stricter model must also be permitted by the weaker one.

\section{Lean 4 and its Meta-Programming Framework}
This subsection introduces the nessary knowleages for those who never used the Lean 4 bofore, we only introduce the most important parts used in this project, we don't introduce complex concepts like monads.
Before we introduce the concepts about the Lean 4, we need to talk about why we choose Lean 4 as our tool in this project.

1. It is an interactive proof assistant (ITP). Most of the proofs for the memory models use pen-and-paper, such as (Podkopaev et al., 2019; Alglave et al., 2014; Batty et al., 2011; Doko and Vafeiadis, 2016). They give theoretical contributions but may have problems, such as they finding errors in existing models (Chakraborty and Vafeiadis, 2016, 2017) and repairing those models (Lahav et al., 2017; Chapter 3). By using the ITP, we could increase the reliability of the proof by avoiding missing edge cases, because the proofs can be checked in a mechanized way by a computer.

2.  Lean 4 is fully extensible: users can modify and extend the parser, elaborator, tactics, decision
procedures, pretty printer, and code generator. (cite: https://lean-lang.org/papers/lean4.pdf), for this project, we separate the paper-and-pen (Linux kernel model).

3. Another option is we use the SMT solver.

The Lean 4 is a relatively new proof assistant currently under the active developement, it's featured by the dependent type. The most powerful part of the lean 4 is its general purpose programming and the metaprogramming framework, that's also why we choose Lean 4 as our mechnizazing tool. (Why it's metaprogramming framework is powerful, compared to what)



\subsubsection{Defintations in Lean 4}
We need to constraint our design, because the Cat langauge itself is a big project and implemented by several years, we can't implement all the parts within a master thesis work, so the target for this project is to try to support the proof of the LKMM model.

The Cat language model is formalled
\label{subsubsection:Subsubsection of section 1}
\lipsum[1-2]

The most foundamental type in Cat langauge is the Event, so to define it correctly is quite important as the starting point. The event is an instruction that reads or writes a value on an address


The cat language is implementated in OCaml, and some of the operations are implemented by the computation, for example, the OCaml implementation of co is use a function called cross, and it's used to get pick up the set from a set of set in order, and we can express it in Lean 4 because Lean 4 is also a general purpose programming languge, but we need to prove lots of proofs for the computation and it's not the main idea of our project. So there is a design choice, we will add the std.cat and cross.cat these basic cat as part of the assumptions in our model instead of parsing all the individual cat files.
For example, the [introduce the computation of co]. In our model, we just assume the computation is finished, and they gave us the results of the computation: co is linearization order.



